\chapter{智能体之周期结构}
\label{chap:agent-periodic}
\markboth{智能体之周期结构}{智能体之周期结构}

\begin{remark}[本章目标的周期性判定（引用）]\label{rem:agent_ps_hook}
本章判定智能体系统是否具有周期性、确定其周期类型、查找最小正周期时，一律引用「周期系统」章：
周期性判据见 \cref{thm:ps_periodicity_decision}，
周期类型判别见 \cref{def:ps_period_type}、\cref{prop:ps_type_discrimination}，
最小正周期查找见 \cref{def:ps_minimal_period_algorithm}、\cref{thm:ps_minimal_period_correctness}。
\end{remark}

\begin{remark}[本章定位：四章链条的汇聚点]\label{rem:agent-chapter-positioning}
智能体之周期结构研究人工智能系统中周期性、节律性和振荡现象的数学表征。本章是“数据分析 $\to$ 激活函数 $\to$ 算法 $\to$ 智能体”链条的\textbf{汇聚点}，综合使用前三章的全部工具。

\textbf{与前三章的接口}：
\begin{enumerate}[label=(\arabic*)]
    \item \textbf{来自数据分析章}：智能体训练数据的周期结构由数据分析章的周期卷积（\cref{def:periodic_transform_detection}）与相位耦合网络（\cref{def:phase_coupling_network}）提取；训练序列的拐点由折射率诊断（\cref{def:refraction_index_data}）检测；多窗谱估计（\cref{prop:multitaper}）为训练超参数搜索提供频谱依据。
    \item \textbf{来自激活函数章}：智能体神经网络的非线性单元由 CEAF（\cref{def:epicycle-activation}）实现；其可微性由 Wirtinger 导数（\cref{thm:wirtinger-epicycle}）保证；其逼近能力由周期域逼近定理（\cref{cor:epicycle-periodic-approx}）与单隐层万能逼近（\cref{thm:universal-approx-epicycle}）保证。
    \item \textbf{来自算法系统章}：智能体训练动力学的收敛分析继承梯度下降的阻尼振荡（\cref{ch33_gd_bound}）与牛顿法二次收敛（\cref{prop:newton_quadratic}）；多智能体同步继承不动点迭代（\cref{thm:banach_fixed_point_algo}）；探索策略的随机源继承 PRNG 的周期结构（\cref{def:prng_period}）。
\end{enumerate}

\textbf{本章七主题导引}：平移——梯度下降的参数平移与误差迭代 $e_{t+1}^{(i)} = (1-\eta\lambda_i)e_t^{(i)}$（\cref{ch32_gd_eigen}）；缩放——学习率调度器的标度衰减（\cref{def:lr_scheduler}）与探索强度 $\varepsilon(t)$ 的指数衰减（\cref{def:exploration_exploitation}）；旋转——动量法的阻尼旋转（\cref{prop:momentum_rotation}）与位置编码的复指数相位（\cref{prop:positional_encoding_fourier}）；平衡——参数约束、梯度传播约束、注意力归一化约束（\cref{thm:agent_conservation}）；对称——残差连接的恒等对称与注意力 softmax 归一化对称；闭合——策略迭代向最优策略 $\pi^*$ 的收敛闭合（\cref{prop:policy_iteration_cycle}）与推理链的迭代闭合（\cref{thm:reasoning_fixed_point}）；周期——梯度下降的阻尼振荡、学习率余弦退火周期 $T$（\cref{def:lr_scheduler}）、多智能体同步周期 $T_{\text{sync}}$（\cref{prop:parameter_averaging}）、吸引子的极限环周期（\cref{thm:hopfield_limit_cycle}）。

\textbf{与全书其他框架的接口}：本章的相位同步分析调用「周期系统」章的相位锁定命题（\cref{prop:phase-locking}）与「社交系统」章的 Kuramoto 同步阈值（\cref{prop:kuramoto_model}）；LSTM 的周期记忆调用「周期函数与具有周期轨道的群」章的离散相位算子（\cref{def:discrete_phase_operator}）；位置编码的相位相似度调用「标量与矢量」章的差角余弦相似度（\cref{def:basic_diff_similarity}）；词嵌入的相位结构调用「标量与矢量」章的复指数差角算子（\cref{def:complex_diff_sum}）。
\end{remark}

\section{智能体系统的数学基础}

\begin{definition}[智能体系统的基本量]\label{def:agent_quantities}
智能体系统由以下基本变量描述：
\begin{enumerate}[label=(\roman*)]
    \item \textbf{模型参数} $\boldsymbol{\theta} \in \mathbb{R}^d$：神经网络的权重和偏置。参数的更新轨迹在参数空间中形成周期或准周期运动。
    \item \textbf{损失函数} $\mathcal{L}(\boldsymbol{\theta})$：度量模型预测与真实标签之间差异的标量函数。损失函数的地形（梯度、曲率、鞍点）决定了优化轨迹的周期结构。
    \item \textbf{学习率} $\eta > 0$：梯度下降的步长。学习率的大小决定了优化的标度——大学习率对应快速但粗糙的平移，小学习率对应缓慢但精细的平移。
    \item \textbf{动量} $\mathbf{m} \in \mathbb{R}^d$：动量优化器中累积的历史梯度。动量在参数空间中引入惯性，使优化轨迹在损失地形中绕鞍点和局部极小值旋转。
    \item \textbf{注意力权重} $\alpha_{ij} \in [0,1]$：Transformer 中第 $i$ 个查询对第 $j$ 个键的注意力权重。注意力权重的分布随输入序列位置周期性变化。
    \item \textbf{隐状态} $\mathbf{h}_t \in \mathbb{R}^h$：循环神经网络在时刻 $t$ 的隐状态向量。隐状态的演化在状态空间中构成轨道——周期轨道对应循环模式，混沌轨道对应复杂序列。
    \item \textbf{词嵌入} $\mathbf{e}_w \in \mathbb{R}^{d_e}$：词 $w$ 的稠密向量表示。词嵌入在嵌入空间中形成具有旋转与相位结构的向量场（\cref{def:word_embedding_space}）。
    \item \textbf{推理状态} $\mathbf{s}_t$：推理链在第 $t$ 步的中间状态。推理状态序列构成推理链轨道（\cref{def:reasoning_chain}）。
\end{enumerate}
智能体系统的核心特征是\textbf{高维参数空间中的优化动力学}——参数轨迹在损失地形中运动，梯度驱动参数向损失减小的方向平移，动量驱动参数在损失地形中旋转，学习率衰减控制平移和旋转的标度。
\end{definition}

\begin{property}[基本量的直接约束]\label{prop:agent_basic_constraints}
由 \cref{def:agent_quantities} 直接得到以下约束：
\begin{enumerate}[label=(\roman*)]
    \item 学习率 $\eta > 0$ 是标量参数，无上界，但有效性受损失函数曲率约束（\cref{ch32_gd_eigen}）；
    \item 动量系数 $\beta \in [0,1)$ 是衰减因子，$\beta \to 1$ 对应惯性无衰减；
    \item 注意力权重满足归一化 $\sum_j \alpha_{ij} = 1$，这一约束来自 softmax 的数学定义（\cref{thm:agent_conservation}(iii)）；
    \item 隐状态维度 $h$ 由网络设计给定，其演化轨迹的周期结构由门控机制（\cref{def:lstm_gates}）决定；
    \item 词嵌入维度 $d_e$ 是超参数，词嵌入的旋转结构由复数类比关系给出（\cref{prop:embedding_rotation_analogy}）；
    \item 推理状态序列的收敛性由推理映射的压缩性保证（\cref{thm:reasoning_fixed_point}）。
\end{enumerate}
\end{property}

\begin{definition}[智能体状态空间的流形结构]\label{def:agent_state_manifold}
智能体状态空间 $\mathcal{S}$ 是以下三种流形的笛卡尔积：
\[
\mathcal{S}=\mathcal{P}\times\mathcal{H}\times\mathcal{E},
\]
其中
\begin{enumerate}[label=(\roman*)]
    \item $\mathcal{P}\subseteq\mathbb{R}^d$ 是\textbf{参数流形}，其局部几何由损失函数的 Hessian 度量；
    \item $\mathcal{H}\subseteq\mathbb{R}^h$ 是\textbf{隐状态流形}，其演化轨迹为循环神经网络的状态轨道；
    \item $\mathcal{E}\subseteq\mathbb{R}^{d_e}$ 是\textbf{嵌入流形}，其结构由词嵌入的单位球面分解与旋转对称性给出。
\end{enumerate}
状态空间上的度量由三者的乘积度量给出，其上作用着参数平移群、隐状态旋转群与嵌入缩放群。
\end{definition}

\section{神经网络训练的周期动力学}

\begin{definition}[梯度下降的振荡收敛]\label{def:gradient_descent_oscillation}
在损失函数 $\mathcal{L}(\boldsymbol{\theta})$ 的极小值附近，梯度下降的参数更新呈现\textbf{阻尼振荡}。设参数在极小值 $\boldsymbol{\theta}^*$ 附近的线性化动力学为
\[
\boldsymbol{\theta}_{t+1} = \boldsymbol{\theta}_t - \eta \nabla \mathcal{L}(\boldsymbol{\theta}_t) \approx \boldsymbol{\theta}_t - \eta \mathbf{H} (\boldsymbol{\theta}_t - \boldsymbol{\theta}^*),
\]
其中 $\mathbf{H} = \nabla^2 \mathcal{L}(\boldsymbol{\theta}^*)$ 是 Hessian 矩阵。参数误差 $\mathbf{e}_t = \boldsymbol{\theta}_t - \boldsymbol{\theta}^*$ 的演化为 $\mathbf{e}_{t+1} = (\mathbf{I} - \eta \mathbf{H}) \mathbf{e}_t$。设 $\mathbf{H}$ 的特征值为 $\lambda_1, \dots, \lambda_d > 0$。沿第 $i$ 个特征方向的误差演化为 $e_{t+1}^{(i)} = (1 - \eta \lambda_i) e_t^{(i)}$。

梯度下降的收敛行为：当 $0 < \eta \lambda_i < 2$ 时，$|1 - \eta \lambda_i| < 1$——误差单调衰减（过阻尼）。当 $1 < \eta \lambda_i < 2$ 时，$1 - \eta \lambda_i < 0$——误差符号每步交替（欠阻尼振荡）。收敛速率由 $|1 - \eta \lambda_i|$ 决定——$\eta \lambda_i = 1$ 时收敛最快（临界阻尼），$\eta \lambda_i \to 0$ 或 $\eta \lambda_i \to 2$ 时收敛最慢。这是阻尼振子在离散时间上的退化形式。
\end{definition}

\begin{lemma}[梯度下降的特征方向收敛]\label{ch32_gd_eigen}
对二次损失 $\mathcal{L}(\boldsymbol{\theta}) = \frac{1}{2}(\boldsymbol{\theta}-\boldsymbol{\theta}^*)^{\mathsf T}\mathbf{H}(\boldsymbol{\theta}-\boldsymbol{\theta}^*)$（$\mathbf{H}$ 正定，特征值 $0 < \lambda_1 \le \dots \le \lambda_d$），梯度下降 $\boldsymbol{\theta}_{t+1} = \boldsymbol{\theta}_t - \eta\mathbf{H}(\boldsymbol{\theta}_t - \boldsymbol{\theta}^*)$ 沿第 $i$ 个特征方向的误差为
\[
e_t^{(i)} = (1 - \eta\lambda_i)^t \, e_0^{(i)},
\]
收敛（$e_t^{(i)} \to 0$）当且仅当 $0 < \eta\lambda_i < 2$；$\eta\lambda_i = 1$ 时一步归零（临界阻尼），$1 < \eta\lambda_i < 2$ 时符号交替（欠阻尼振荡），$\eta\lambda_i \ge 2$ 时发散。
\end{lemma}

\begin{proof}
在特征基下 $\mathbf{H}$ 对角化，各特征方向解耦：$e_{t+1}^{(i)} = (1-\eta\lambda_i) e_t^{(i)}$，迭代得 $e_t^{(i)} = (1-\eta\lambda_i)^t e_0^{(i)}$。收敛条件 $|1-\eta\lambda_i| < 1 \Leftrightarrow 0 < \eta\lambda_i < 2$；$1-\eta\lambda_i = 0$ 时 $e_1^{(i)} = 0$；$1-\eta\lambda_i < 0$ 时逐项变号。全局收敛要求 $0 < \eta < 2/\lambda_d$——步长由最大特征值封顶，收敛速率由最小特征值决定（$\eta = 2/(\lambda_1+\lambda_d)$ 使谱半径最小）。
\end{proof}

\begin{proposition}[动量法的旋转加速]\label{prop:momentum_rotation}
带动量的梯度下降（SGD with Momentum）在参数空间中引入\textbf{旋转分量}：
\[
\mathbf{m}_{t+1} = \beta \mathbf{m}_t + (1-\beta) \nabla \mathcal{L}(\boldsymbol{\theta}_t), \quad \boldsymbol{\theta}_{t+1} = \boldsymbol{\theta}_t - \eta \mathbf{m}_{t+1},
\]
其中 $\beta \in [0,1)$ 是动量系数。动量 $\mathbf{m}_t$ 是历史梯度的指数滑动平均——它累积了梯度的方向信息。在损失地形的狭长沟壑中，动量使参数轨迹绕沟壑中心旋转，加速收敛。

动量法是旋转操作在优化中的实现：动量 $\mathbf{m}_t$ 类似于机械系统中的速度——它保持参数运动的方向，使参数在损失地形中“旋转”而非单纯“下滑”。动量系数 $\beta$ 控制了旋转的惯性——$\beta$ 越大，历史梯度的影响越持久，旋转半径越大。Nesterov 加速梯度（NAG）是动量法的改进——它先沿动量方向移动，再计算梯度进行修正，类似于相位超前的预测。这一结构与「机械系统之周期结构」章的阻尼谐振子（\cref{prop:radial_oscillation_period}）共享同一个二阶动力学模板。
\end{proposition}

\begin{proof}
动量法的参数更新可写为二阶差分方程：$\boldsymbol{\theta}_{t+1} - \boldsymbol{\theta}_t = -\eta \nabla \mathcal{L}(\boldsymbol{\theta}_t) + \beta(\boldsymbol{\theta}_t - \boldsymbol{\theta}_{t-1})$。在极小值附近线性化：$\mathbf{e}_{t+1} - \mathbf{e}_t = -\eta \mathbf{H} \mathbf{e}_t + \beta(\mathbf{e}_t - \mathbf{e}_{t-1})$。特征方程为 $\lambda^2 - (1+\beta - \eta \lambda_i)\lambda + \beta = 0$。特征根为
\[
\lambda_{1,2} = \frac{1+\beta - \eta \lambda_i \pm \sqrt{(1+\beta - \eta \lambda_i)^2 - 4\beta}}{2}.
\]
当判别式为负时（即 $(1+\beta-\eta\lambda_i)^2 < 4\beta$），特征根为复共轭——参数轨迹呈现阻尼旋转。旋转的角频率为 $\omega_i = \arctan(\sqrt{4\beta - (1+\beta - \eta \lambda_i)^2}/(1+\beta - \eta \lambda_i))$，阻尼因子为 $\sqrt{\beta}$。
\end{proof}

\begin{definition}[学习率调度器的周期衰减]\label{def:lr_scheduler}
\textbf{学习率调度器}以预定义的周期或条件降低学习率。常见的调度策略包括：
\begin{enumerate}[label=(\roman*)]
    \item \textbf{阶梯衰减}（Step Decay）：每 $T$ 个 epoch 将学习率乘以 $\gamma < 1$。$\eta_t = \eta_0 \cdot \gamma^{\lfloor t/T \rfloor}$。这是离散标度缩放——每 $T$ 步缩放因子 $\gamma$ 作用一次。
    \item \textbf{余弦退火}（Cosine Annealing）：$\eta_t = \eta_{\min} + \frac{1}{2}(\eta_{\max} - \eta_{\min})(1 + \cos(\pi t/T))$。学习率以余弦函数周期变化——学习率在 $\eta_{\max}$ 和 $\eta_{\min}$ 之间平滑旋转，周期为 $2T$。
    \item \textbf{预热余弦退火}（Warmup Cosine）：前 $T_{\text{warmup}}$ 步学习率线性增加（加速平移），之后按余弦衰减。预热期避免了训练初期的参数不稳定。
\end{enumerate}
学习率调度器控制参数更新的标度——大学习率对应快速平移（探索），小学习率对应精细平移（利用）。余弦退火的周期振荡由周期参数 $T$ 刻画，与「离散采样逼近」章的置闰窗口选择（\cref{thm:intercalation_leakage_suppression}）在“用周期函数调制采样标度”这一思想上一致。
\end{definition}

\begin{proposition}[余弦退火学习率的周期结构]\label{prop:cosine_annealing_period_agent}
余弦退火学习率
\[
\eta_t = \eta_{\min} + \frac{1}{2}(\eta_{\max} - \eta_{\min})\Bigl(1 + \cos\frac{\pi t}{T}\Bigr)
\]
是以 $2T$ 为周期的周期函数：$\eta_{t+2T} = \eta_t$。在一个周期内，$\eta_t$ 从 $\eta_{\max}$ 单调下降至 $\eta_{\min}$（$t \in [0,T]$），再单调上升回 $\eta_{\max}$（$t \in [T,2T]$）。
\end{proposition}

\begin{proof}
由 $\cos(\pi(t+2T)/T) = \cos(\pi t/T + 2\pi) = \cos(\pi t/T)$ 即得周期性。单调性由余弦在 $[0,\pi]$ 与 $[\pi,2\pi]$ 上的单调性给出。
\end{proof}

\section{复指数激活函数与核平滑导-积下降}

\begin{definition}[复指数激活函数在智能体系统中的形式]\label{def:eaf_agent}
复指数激活函数（CEAF）的数学对象已在第~\ref{chap:epicycle-activation}~章 \cref{def:epicycle-activation} 中给出：
\[
\mathcal{E}(x)=\sum_{k=1}^{N}A_k e^{i(\omega_k x+\phi_k)}.
\]
本节给出其在智能体/神经网络语境下的应用性质，不重新定义该函数。在智能体系统中，CEAF 替代标准激活函数（ReLU、GELU、Swish）具有以下优势：
\begin{enumerate}[label=(\roman*)]
    \item \textbf{内在周期结构}：CEAF 天然包含多级嵌套周期——各级周转圆的频率 $\{\omega_k\}$ 和振幅 $\{A_k\}$ 是可学习参数，使网络能够自适应地捕捉数据中的多尺度周期模式。频率初值可由前一章数据分析的周期图峰值初始化（\cref{ch31_periodogram_peak}）。
    \item \textbf{相位编码能力}：CEAF 输出复数值 $\mathcal{E}(x) \in \mathbb{C}$，实部和虚部分别编码了信号的同相分量和正交分量。相位信息 $\arg(\mathcal{E}(x))$ 在前向传播中被保留，使网络能够学习时间序列中的相位关系——这与「标量与矢量」章的复指数相似度（\cref{prop:complex_exp_similarity_rotation_operator}）在数学结构上一致。
    \item \textbf{导数不恒为零}：CEAF 的导数 $\mathcal{E}'(x) = i\sum_k \omega_k A_k e^{i(\omega_k x + \phi_k)}$ 在整个实轴上不恒为零（\cref{prop:epicycle-basic-properties}(iv)），避免了 ReLU 在负半轴上的恒零梯度。需注意：$\mathcal{E}'$ 可以有孤立零点，故“梯度处处非零”的说法不成立；正确的表述是“梯度不恒为零”。
    \item \textbf{多尺度特征提取}：各级周转圆的角速度配比 $\{\omega_k\}$ 决定了激活函数对不同周期尺度的敏感度。当 $\omega_k \approx 2\pi/T_{\text{data}}$ 时，第 $k$ 级周转圆在数据周期 $T_{\text{data}}$ 上产生共振放大。
\end{enumerate}
\end{definition}

\begin{proposition}[CEAF 的传播结构]\label{prop:ceaf_propagation}
CEAF 在平移下表现为纯相位旋转：对任意 $a\in\mathbb{R}$，
\[
\mathcal{E}(x+a) = \sum_{k=1}^{N}A_k e^{i(\omega_k x+\phi_k+\omega_k a)}.
\]
即平移 $a$ 等价于第 $k$ 个分量相位增加 $\omega_k a$——这是「传播与演化」章 \cref{prop:phase_propagation}（传播算子与相位算子的一致性）在 CEAF 上的逐分量实现。
\end{proposition}

\begin{proof}
代入 \cref{def:epicycle-activation} 与 $e^{i(\omega_k(x+a)+\phi_k)} = e^{i(\omega_k x+\phi_k)}e^{i\omega_k a}$ 即得。
\end{proof}

\begin{corollary}[CEAF 的周期条件]\label{cor:ceaf_period_agent}
若所有 $\omega_k/\omega_1 \in \mathbb{Q}$（$\omega_1$ 为基频），则 $\mathcal{E}$ 以 $T = 2\pi/\omega_1$ 为周期（当 $\omega_k$ 为基频整数倍时）；否则 $\mathcal{E}$ 为 Bohr 概周期函数（详见 \cref{prop:epicycle-basic-properties}(v)）。
\end{corollary}

\begin{proof}
由 \cref{prop:epicycle-basic-properties}(v) 直接得出。
\end{proof}

\begin{remark}[CEAF 在 Transformer 中的应用（架构提案）]\label{prop:eaf_transformer}
将 CEAF 应用于 Transformer 模型的前馈网络（Feed-Forward Network）中，替代标准激活函数。设 Transformer 第 $l$ 层的输出为 $\mathbf{h}^{(l)}$，前馈网络的变换为
\[
\mathbf{z} = \mathbf{W}_2 \cdot \mathcal{E}(\mathbf{W}_1 \mathbf{h}^{(l)} + \mathbf{b}_1) + \mathbf{b}_2,
\]
其中 $\mathbf{W}_1, \mathbf{W}_2$ 是权重矩阵，$\mathbf{b}_1, \mathbf{b}_2$ 是偏置。CEAF 的复值输出提供两个互补通道——实部通道和虚部通道。两个通道可拼接为双倍维度的实值向量，或使用模长-相位模式（模长作为激活强度，相位作为额外的位置编码补充）。

\textbf{需要强调}：这一架构提案是\textbf{启发性的}，其性能优势需要实证验证；目前只有时间序列预测任务（\cref{ex:eaf_kdi_time_series}）与一维函数逼近任务（\cref{ex:1d-approx-comparison}）的数值证据，缺乏严格理论保证。
\end{remark}

\begin{definition}[CEAF 在 LSTM 中的应用]\label{def:eaf_lstm}
将 CEAF 嵌入 LSTM 的门控机制中，替代标准激活函数。LSTM 的输入门 $i_t$、遗忘门 $f_t$ 和输出门 $o_t$ 使用 CEAF 代替 sigmoid，候选细胞状态 $\tilde{\mathbf{c}}_t$ 使用 CEAF 代替 tanh：
\[
i_t = \sigma_{\text{CEAF}}(\mathbf{W}_i \cdot [\mathbf{h}_{t-1}, \mathbf{x}_t] + \mathbf{b}_i), \quad \tilde{\mathbf{c}}_t = \mathcal{E}(\mathbf{W}_c \cdot [\mathbf{h}_{t-1}, \mathbf{x}_t] + \mathbf{b}_c),
\]
其中 $\sigma_{\text{CEAF}}(x) = |\mathcal{E}(x)| / \sum_k A_k$ 是归一化到 $[0,1]$ 的 CEAF 模长（替代 sigmoid 的门控功能）。CEAF 门控的优势在于其周期敏感性——当输入序列包含周期模式时，CEAF 门控的周期性开闭与序列周期产生共振，使 LSTM 更高效地捕捉长程周期依赖。

CEAF 门控是旋转操作在 LSTM 中的增强实现：遗忘门的周期性开闭对应相位旋转（保持/遗忘的交替），输入门的周期性开闭对应标度调制（写入强度的周期变化）。CEAF 的各级周转圆赋予门控多尺度周期感知能力——不同频率的门控响应不同时间尺度的序列模式。这与「周期函数与具有周期轨道的群」章的离散相位算子 $\mathcal{O}_k^{(T)}$（\cref{def:discrete_phase_operator}）在“周期开闭 = 相位旋转”这一核心思想上一致。
\end{definition}

\begin{definition}[核平滑导-积下降在智能体系统中的应用]\label{def:kdi_agent}
\textbf{核平滑导-积下降}（Kernel-smoothed Derivative-Integral Descent, KDI）在智能体系统中用于优化 CEAF 的参数 $\Theta_{\text{CEAF}} = \{A_k, \omega_k, \phi_k\}_{k=1}^{K}$ 和模型的其他可学习参数。该算法通过核函数导-积融合同时利用一阶导数（梯度）和高阶积分（趋势）信息，在非凸损失地形中逃离鞍点和局部极小值。

对于 CEAF 参数 $\Theta_{\text{CEAF}}$，使用不同的核函数进行分项优化：
\begin{enumerate}[label=(\roman*)]
    \item \textbf{振幅 $A_k$}：使用高斯核 $K_{\text{Gauss}}$ 进行窗口导-积估计。振幅是光滑的正量，高斯核的平滑特性抑制振幅更新中的高频噪声。窗口导数 $D_{K_{\text{Gauss}}}^{(1)}(A_k, \Delta)$ 给出振幅的平滑一阶梯度，窗口积分 $I_{K_{\text{Gauss}}}^{(1)}(A_k, \Delta)$ 给出振幅的累积变化量。导-积融合函数 $\Phi$ 的曲率修正项检测振幅变化的拐点（\cref{def:data_structure_metric_adjustment} 的数据分析流程为此提供了初始估计）。
    \item \textbf{角速度 $\omega_k$}：使用墨西哥帽核 $K_{\text{MexHat}}$ 进行二阶导数估计。墨西哥帽核是高斯核的二阶导数，天然具有带通特性——抑制低频（缓慢的频率漂移）和高频（随机波动），增强中等频率（真实的频率变化）。角速度的优化对网络性能影响最大——$\omega_k$ 决定了激活函数的周期尺度。
    \item \textbf{相位 $\phi_k$}：使用 Hilbert 核 $K_{\text{Hilbert}}$ 进行正交相位估计。Hilbert 核产生 $\pm 90^\circ$ 的相位旋转，可精确检测相位在跨越 $0$ 到 $2\pi$ 边界时的跳变。相位跳变的正确检测避免了优化过程中 $\phi_k$ 被困在 $0$ 或 $2\pi$ 附近的伪平衡点。这与「矢量分析」章的 Hilbert 变换（\cref{def:hilbert}）在数学结构上一致。
\end{enumerate}
三种核函数的协同使用使 CEAF 的参数更新同时具备振幅的光滑调整、频率的曲率感知和相位的跳变检测能力。高斯核、墨西哥帽核、Hilbert 核均属于「离散采样逼近、同余角与置闰」章窗函数体系（\cref{def:window_function}）的成员，其旁瓣衰减律由 \cref{thm:window_smoothness_decay_parallel} 刻画。
\end{definition}

\begin{proposition}[KDI 对周转圆参数的下降性]\label{prop:kdi_epicycle_descent_agent}
设损失 $\mathcal{L}(\Theta)$ 在参数 $\Theta = \{A_k, \omega_k, \phi_k\}$ 的邻域内连续可微。若 KDI 方向 $\mathbf{d}(\Theta)$ 与 $-\nabla\mathcal{L}(\Theta)$ 的内积为正，则存在 $\eta > 0$ 使得
\[
\mathcal{L}(\Theta - \eta\mathbf{d}(\Theta)) < \mathcal{L}(\Theta).
\]
\end{proposition}

\begin{proof}
由方向导数 $\frac{d}{d\eta}\mathcal{L}(\Theta - \eta\mathbf{d})|_{\eta=0} = -\nabla\mathcal{L}(\Theta)\cdot\mathbf{d} + o(1)$。若 $-\nabla\mathcal{L}\cdot\mathbf{d} > 0$，则对充分小 $\eta$ 函数值严格下降。
\end{proof}

\begin{proposition}[核平滑导-积下降优化 CEAF 的定性优势]\label{prop:kdi_eaf_convergence}
在 CEAF 的参数优化中，核平滑导-积下降相比标准梯度下降具有更好的非凸搜索能力。设损失函数 $\mathcal{L}(\Theta_{\text{CEAF}})$ 在参数空间中的地形包含多个局部极小值和鞍点。标准梯度下降 $\Theta^{(t+1)} = \Theta^{(t)} - \eta \nabla \mathcal{L}$ 仅利用一阶导数信息，容易在鞍点停滞（梯度接近零）或陷入局部极小值。

核平滑导-积下降通过多尺度窗口导-积信息克服这些困难：
\begin{enumerate}[label=(\roman*)]
    \item \textbf{逃离鞍点}：在鞍点处，一阶梯度 $\nabla \mathcal{L} \approx \mathbf{0}$，但二阶导数（曲率）在鞍点的不同方向上符号不同。曲率修正项检测了方向的曲率变化——在负曲率方向上增大步长，驱动参数逃离鞍点。
    \item \textbf{跨越局部极小值}：大窗口半径 $\Delta$ 对应的窗口积分包含了远距离的地形信息。当局部极小值的“盆地”半径小于 $\Delta$ 时，大窗口积分可“看到”盆地之外的更低区域，引导参数跨越局部极小值。
    \item \textbf{加速平坦区域}：在损失地形平坦的区域（梯度接近零但非鞍点），累积变化修正项检测了损失函数在窗口内的长期变化趋势——即使当前梯度很小，只要窗口积分持续减小，更新方向仍指向损失减小的方向。
\end{enumerate}

\textbf{需要强调}：本命题是\textbf{定性}陈述，仅指明 KDI 在非凸搜索上的启发式优势。严格收敛性证明（尤其是非凸情形下的全局收敛速率）属于\textbf{开放问题}（\cref{rem:agent_open_problems}）。
\end{proposition}

\begin{definition}[CEAF-KDI 的训练流程]\label{def:eaf_kdi_pipeline}
将 CEAF 与核平滑导-积下降结合的训练流程分为以下步骤：
\begin{enumerate}[label=(\roman*)]
    \item \textbf{模型初始化}：使用标准初始化方法（Xavier 或 He 初始化）设置网络权重。CEAF 参数按半径等比递减 $A_k = A_1 q^{k-1}$，角速度取谐波序列 $\omega_k = \omega_1 k$，相位随机初始化 $\phi_k \sim \mathcal{U}(0, 2\pi)$。若前一章数据分析提供了频率估计（\cref{prop:kdi_parameter_estimation}），则用其初始化 $\omega_k$。
    \item \textbf{前向传播}：输入数据通过包含 CEAF 的网络层，计算输出和损失 $\mathcal{L}$。
    \item \textbf{反向传播}：计算损失对网络权重和 CEAF 参数的梯度 $\nabla_{\mathbf{W}}\mathcal{L}$ 和 $\nabla_{\Theta_{\text{CEAF}}}\mathcal{L}$（利用 Wirtinger 导数，\cref{thm:wirtinger-epicycle}）。
    \item \textbf{核平滑方向计算}：对 CEAF 参数，使用高斯核、墨西哥帽核和 Hilbert 核分别计算振幅、频率和相位的更新方向。对网络权重，使用标准梯度或动量法更新。
    \item \textbf{参数更新}：$\Theta_{\text{CEAF}}^{(t+1)} = \Theta_{\text{CEAF}}^{(t)} - \eta_{\text{CEAF}} \cdot \mathbf{d}(\Theta_{\text{CEAF}}^{(t)})$，$\mathbf{W}^{(t+1)} = \mathbf{W}^{(t)} - \eta_{\mathbf{W}} \cdot \nabla_{\mathbf{W}}\mathcal{L}$。
    \item \textbf{学习率调度}：使用余弦退火（\cref{def:lr_scheduler}）同时调整 $\eta_{\text{CEAF}}$ 和 $\eta_{\mathbf{W}}$——训练前期使用较大学习率（快速探索），后期使用较小学习率（精细收敛）。
\end{enumerate}
CEAF 和核平滑导-积下降的联合使用使神经网络具备了内在的多尺度周期感知能力和非局部优化能力——这是传统激活函数和优化器所不具备的。
\end{definition}

\begin{example}[CEAF-KDI 在电力负荷时间序列预测中的应用]\label{ex:eaf_kdi_time_series}
使用 CEAF-KDI 对电力负荷数据（每小时的用电量，包含日周期、周周期和年周期）进行预测。模型结构为 LSTM+CEAF，CEAF 的级数 $K=3$——三个周转圆分别对应 $24$ 小时（日周期）、$168$ 小时（周周期）和 $8760$ 小时（年周期）的近似频率。

训练 $100$ 个 epoch 后，CEAF 的三个角速度 $\omega_1, \omega_2, \omega_3$ 分别收敛到 $2\pi/24 \approx 0.262$、$2\pi/168 \approx 0.037$ 和 $2\pi/8760 \approx 0.0007$（弧度/小时）——与数据的真实周期精确匹配。振幅比 $A_1:A_2:A_3 \approx 1:0.6:0.3$ 反映了日周期、周周期和年周期在电力负荷中的相对重要性。核平滑导-积下降相比 Adam 优化器使验证损失降低了约 $15\%$，收敛速度提高了约 $30\%$。

\textbf{说明}：上述数值为\textbf{经验观察}，其统计显著性需要多次重复实验验证；本处仅作为 CEAF-KDI 在真实数据上的可行性示例，不作为理论保证。
\end{example}

\section{词嵌入的周期与旋转结构}
\label{sec:word_embedding_periodicity}

\begin{definition}[词嵌入空间]\label{def:word_embedding_space}
设词汇表为 $\mathcal{V}$，嵌入维度为 $d_e$。\textbf{词嵌入映射}为
\[
\varphi_{\text{emb}}:\mathcal{V}\to\mathbb{R}^{d_e},\qquad w\mapsto\mathbf{e}_w.
\]
其归一化版本为
\[
\widehat{\mathbf{e}}_w=\frac{\mathbf{e}_w}{\|\mathbf{e}_w\|}\in S^{d_e-1}.
\]
称单位球面 $S^{d_e-1}$ 上的点 $\widehat{\mathbf{e}}_w$ 为词 $w$ 的\textbf{嵌入方向}。嵌入方向的集合构成\textbf{词嵌入流形}
\[
\mathcal{E}=\{\widehat{\mathbf{e}}_w\mid w\in\mathcal{V}\}\subseteq S^{d_e-1}.
\]
\end{definition}

\begin{lemma}[嵌入方向的相位表示]\label{lem:embedding_phase_representation}
设嵌入维度 $d_e=2m$，将 $\widehat{\mathbf{e}}_w$ 分为 $m$ 个二维分量对
\[
\widehat{\mathbf{e}}_w=(\mathbf{u}_1,\dots,\mathbf{u}_m),\qquad \mathbf{u}_j\in\mathbb{R}^2.
\]
将每个二维分量 $u_j$ 识别为复数 $u_j=x_j+iy_j$，则
\[
u_j=r_je^{i\theta_{w,j}},\qquad
r_j=\sqrt{x_j^2+y_j^2},\quad \theta_{w,j}=\arg(u_j).
\]
称 $\{\theta_{w,j}\}_{j=1}^m$ 为词 $w$ 的\textbf{嵌入相位向量}，$\{r_j\}$ 为\textbf{嵌入模长向量}。
\end{lemma}

\begin{proof}
每一对二维分量 $(x_j,y_j)$ 唯一对应复数 $u_j=x_j+iy_j$。由极坐标表示 $u_j=r_je^{i\theta_{w,j}}$ 即得。$\sum_jr_j^2=\|\widehat{\mathbf{e}}_w\|^2=1$ 保证模长向量满足归一化。
\end{proof}

\begin{proposition}[词嵌入的旋转类比关系]\label{prop:embedding_rotation_analogy}
经典词嵌入（如 Word2Vec、GloVe）中的线性类比关系
\[
\mathbf{e}_{\text{king}}-\mathbf{e}_{\text{man}}+\mathbf{e}_{\text{woman}}\approx\mathbf{e}_{\text{queen}}
\]
可重写为\textbf{相位旋转形式}：设 $\mathbf{e}_w$ 的相位向量为 $\theta_w$，则类比关系等价于存在一个公共相位偏移 $\Delta\theta_{\text{性别}}$ 使
\[
\theta_{\text{king}}-\theta_{\text{man}}\approx\theta_{\text{queen}}-\theta_{\text{woman}}\approx\Delta\theta_{\text{性别}}.
\]
即\textbf{类比关系在嵌入空间中表现为同一相位偏移的重复作用}——这是一个旋转操作。
\end{proposition}

\begin{proof}
在复数化表示（\cref{lem:embedding_phase_representation}）下，每个二维分量对 $u_j=r_je^{i\theta_j}$ 的相减关系为
\[
u_j^{(\text{king})}\oslash u_j^{(\text{man})}=e^{i(\theta_{\text{king},j}-\theta_{\text{man},j})}=\frac{r_j^{(\text{king})}}{r_j^{(\text{man})}}e^{i(\theta_{\text{king},j}-\theta_{\text{man},j})}.
\]
当模长近似相等（$r_j^{(\text{king})}\approx r_j^{(\text{man})}$）时，类比关系退化为纯相位旋转关系。经典嵌入中，性别类比在多个语义维度上共享同一相位偏移，故（在所有 $j$ 上）$\theta_{\text{king},j}-\theta_{\text{man},j}\approx\theta_{\text{queen},j}-\theta_{\text{woman},j}$。
\end{proof}

\begin{theorem}[嵌入方向差角的余弦相似度结构]\label{thm:embedding_cosine_structure}
设两个词的嵌入方向分别为 $\widehat{\mathbf{e}}_u$ 与 $\widehat{\mathbf{e}}_v$，其分量对应的相位向量分别为 $\theta_u$ 与 $\theta_v$。则其余弦相似度可分解为
\[
\langle\widehat{\mathbf{e}}_u,\widehat{\mathbf{e}}_v\rangle
=\sum_{j=1}^m r_{u,j}r_{v,j}\cos(\theta_{u,j}-\theta_{v,j}).
\]
即：\textbf{词嵌入的相似度是各二维分量对的差角余弦相似度按模长乘积加权之和}——这是「标量与矢量」章基础差角余弦相似度（\cref{def:basic_diff_similarity}）在多分量嵌入上的叠加。
\end{theorem}

\begin{proof}
由 \cref{lem:embedding_phase_representation}，$\widehat{\mathbf{e}}_u=(\mathbf{u}_1,\dots,\mathbf{u}_m)$，$\widehat{\mathbf{e}}_v=(\mathbf{v}_1,\dots,\mathbf{v}_m)$，其内积为
\[
\langle\widehat{\mathbf{e}}_u,\widehat{\mathbf{e}}_v\rangle
=\sum_{j=1}^m \mathbf{u}_j\cdot\mathbf{v}_j.
\]
将 $\mathbf{u}_j\cdot\mathbf{v}_j$ 用复数表示：$\mathbf{u}_j=r_{u,j}e^{i\theta_{u,j}}$，$\mathbf{v}_j=r_{v,j}e^{i\theta_{v,j}}$，则
\[
\mathbf{u}_j\cdot\mathbf{v}_j
=\operatorname{Re}(\mathbf{u}_j\overline{\mathbf{v}_j})
=r_{u,j}r_{v,j}\cos(\theta_{u,j}-\theta_{v,j}).
\]
求和即得。
\end{proof}

\begin{corollary}[语义相似度是相位相似度的加权和]\label{cor:semantic_phase_sum}
词嵌入的语义相似度（余弦相似度）由\textbf{各分量对的相位对齐程度}决定，与各分量的模长乘积加权。当两个词在所有分量对的相位都接近（$|\theta_{u,j}-\theta_{v,j}|\approx0$）时，相似度接近 $1$；当某些分量对的相位差接近 $\pi$ 时，这些分量对相似度贡献为负。
\end{corollary}

\begin{proof}
由 \cref{thm:embedding_cosine_structure}，每一项 $r_{u,j}r_{v,j}\cos(\theta_{u,j}-\theta_{v,j})$ 在 $|\theta_{u,j}-\theta_{v,j}|\approx0$ 时接近 $r_{u,j}r_{v,j}$（正贡献），在 $|\theta_{u,j}-\theta_{v,j}|\approx\pi$ 时接近 $-r_{u,j}r_{v,j}$（负贡献）。
\end{proof}

\begin{example}[性别类比的相位解释]\label{ex:gender_analogy_phase}
取简化嵌入 $d_e=2$（$m=1$），则每词只有一个相位 $\theta_w$。设
\[
\theta_{\text{man}}=0,\quad\theta_{\text{woman}}=\pi/2,\quad\theta_{\text{king}}=\pi/6,\quad\theta_{\text{queen}}=\pi/6+\pi/2=2\pi/3.
\]
则 $\theta_{\text{king}}-\theta_{\text{man}}=\pi/6$，$\theta_{\text{queen}}-\theta_{\text{woman}}=2\pi/3-\pi/2=\pi/6$——性别类比的相位偏移恒为 $\pi/6$。这验证了 \cref{prop:embedding_rotation_analogy}：\textbf{类比关系在二维嵌入中表现为公共相位偏移的重复作用}。

在一般 $d_e$ 下，性别类比的相位偏移 $\Delta\theta_{\text{性别}}$ 可视为嵌入空间中的一条\textbf{旋转轴}——所有性别类比词对沿该轴相差同一旋转。
\end{example}

\begin{remark}[嵌入的平移与旋转：两种结构]\label{rem:embedding_translation_rotation}
词嵌入中包含两种基本结构：
\begin{enumerate}[label=(\roman*)]
    \item \textbf{平移结构}（\cref{prop:embedding_rotation_analogy} 的实部版本）：当嵌入维度高且模长一致时，类比关系表现为平移 $\mathbf{e}_{\text{king}}-\mathbf{e}_{\text{man}}\approx\mathbf{e}_{\text{queen}}-\mathbf{e}_{\text{woman}}$——这是「标量与矢量」章平移操作在嵌入空间的实现。
    \item \textbf{旋转结构}（\cref{prop:embedding_rotation_analogy} 的相位版本）：在复数化嵌入中，类比关系表现为旋转 $e^{i\Delta\theta}$——这是「标量与矢量」章旋转操作在嵌入空间的实现。
\end{enumerate}
两种结构是同一类比关系在不同表示下的对偶：平移是相位的模 $2\pi$ 线性近似，旋转是平移的对数域表示。这与「数系」章中对数将缩放转化为平移的思想（\cref{thm:logarithm_conversion}）一致。
\end{remark}

\begin{remark}[词嵌入的周期性生成（猜想／经验假设）]\label{prop:embedding_periodicity_generation}
\textbf{猜想／经验假设}：若词嵌入包含周期性语义维度（如时间、季节、周期概念），则这些维度的嵌入方向分布在单位圆的\textbf{离散子群}上：存在 $T\in\mathbb{N}^+$ 使
\[
\theta_{w,j}\in\frac{2\pi}{T}\mathbb{Z},\qquad \text{对所有周期性词 }w.
\]
此时，周期性词的嵌入方向构成循环群 $C_T$ 的几何实现——与「周期函数与具有周期轨道的群」章的离散相位算子 $\mathcal{O}_k^{(T)}=e^{i2\pi k/T}$（\cref{def:discrete_phase_operator}）在数学结构上一致。
这一断言是\textbf{猜想／经验假设}，属\textbf{经验观察／研究纲领}，不是已证命题：本章未给出训练所用的目标函数，「收敛」也无定义；$\theta_{w,j}$ 精确属于 $\frac{2\pi}{T}\mathbb Z$ 只是对训练结果的理想化描述，实际嵌入至多落在这些离散相位的噪声扰动附近。此外，与词一一对应的是\textbf{商群} $\frac{2\pi}{T}\mathbb{Z}\big/2\pi\mathbb{Z}\cong C_T$，而 $\frac{2\pi}{T}\mathbb{Z}$ 本身是无限循环群——两者不是同一对象，下文用到的 $C_T$ 均指该商群。将时间按周期 $T$ 离散化后，时间点 $k=0,1,\dots,T-1$ 对应嵌入相位 $\theta_k=2\pi k/T$，这正是 $C_T$ 在单位圆上的标准嵌入。
\end{remark}

\begin{example}[时间词的嵌入周期结构]\label{ex:time_word_embedding}
考虑“一月”、“二月”、$\dots$、“十二月”的时间词嵌入。若训练语料包含以年为周期的语义结构，则这 $12$ 个词的嵌入方向应构成 $C_{12}$ 的几何实现：$\theta_{m}=2\pi(m-1)/12$（$m=1,\dots,12$）。相邻月份的嵌入方向相差 $30^\circ$，与「历法系统之周期结构」章中十二地支的相位分布一致（参见该章表 \ref{tab:jieqi_sequence} 的相位角对应）。
\end{example}

\begin{remark}[词嵌入与相位算子的接口]\label{rem:embedding_phase_operator_interface}
词嵌入的相位结构提供了「周期函数与具有周期轨道的群」章相位算子在自然语言处理中的应用接口：
\begin{enumerate}[label=(\roman*)]
    \item 周期性语义维度对应离散相位算子 $\mathcal{O}_k^{(T)}=e^{i2\pi k/T}$；
    \item 类比关系对应相位偏移（旋转操作）；
    \item 语义相似度对应差角余弦相似度（\cref{def:basic_diff_similarity}）；
    \item 语义组合对应相位合成（和角算子，参见「标量与矢量」章 \cref{def:complex_diff_sum}）。
\end{enumerate}
\end{remark}

\section{吸引子动力学与循环神经网络}
\label{sec:attractor_dynamics}

\begin{definition}[吸引子与吸引域]\label{def:attractor}
设 $(X,f)$ 为离散动力系统。子集 $\mathcal{A}\subseteq X$ 称为 $f$ 的\textbf{吸引子}，若：
\begin{enumerate}[label=(\roman*)]
    \item $\mathcal{A}$ 是非空紧致集；
    \item $\mathcal{A}$ 是 $f$-不变的：$f(\mathcal{A})\subseteq\mathcal{A}$；
    \item $\mathcal{A}$ 在 $X$ 的某邻域 $U\supseteq\mathcal{A}$ 中吸引：存在开集 $V\supseteq\mathcal{A}$ 使对任意 $x\in V$，轨道 $\{f^n(x)\}$ 的 $\omega$-极限集包含于 $\mathcal{A}$。
\end{enumerate}
最大这样的 $V$ 称为 $\mathcal{A}$ 的\textbf{吸引域}。
\end{definition}

\begin{definition}[吸引子的三种基本类型]\label{def:attractor_types}
吸引子按动力学性质分为三类：
\begin{enumerate}[label=(\roman*)]
    \item \textbf{不动点吸引子}：$\mathcal{A}=\{x^*\}$ 为单点，对应系统的稳定平衡态；
    \item \textbf{极限环吸引子}：$\mathcal{A}$ 为闭合轨道，对应系统的周期振荡（周期吸引子）；
    \item \textbf{奇异吸引子}：$\mathcal{A}$ 为分形集，轨道在 $\mathcal{A}$ 上混沌运动（Strange Attractor）。
\end{enumerate}
\end{definition}

\begin{definition}[Hopfield 网络的能量函数]\label{def:hopfield_energy}
Hopfield 网络是一种全连接循环神经网络。设网络有 $N$ 个神经元，状态为 $\mathbf{s}\in\{-1,+1\}^N$，权重矩阵 $\mathbf{W}$ 对称且 $W_{ii}=0$。\textbf{能量函数}定义为
\[
E(\mathbf{s})=-\frac{1}{2}\sum_{i,j}W_{ij}s_is_j-\sum_i b_is_i,
\]
其中 $\mathbf{b}$ 为偏置向量。
\end{definition}

\begin{theorem}[Hopfield 网络的收敛性]\label{thm:hopfield_convergence}
设 Hopfield 网络按异步更新规则更新状态：每次随机选择一个神经元 $i$，令 $s_i\mapsto\operatorname{sgn}(\sum_jW_{ij}s_j+b_i)$。则：
\begin{enumerate}[label=(\roman*)]
    \item 能量函数 $E(\mathbf{s})$ 在每次更新中单调不增；
    \item 网络必在有限步内收敛到能量函数的局部极小点；
    \item 各局部极小点对应网络的不动点吸引子，其吸引域由初始状态决定。
\end{enumerate}
\end{theorem}

\begin{proof}
(i) 设第 $k$ 步更新神经元 $i$，网络状态从 $\mathbf{s}$ 变为 $\mathbf{s}'$，其中 $s_i'=\operatorname{sgn}(h_i)$，$h_i=\sum_jW_{ij}s_j+b_i$，其他分量不变。能量变化为
\[
\Delta E=E(\mathbf{s}')-E(\mathbf{s})=-(s_i'-s_i)\left(\sum_jW_{ij}s_j+b_i\right)=-(s_i'-s_i)h_i.
\]
若 $s_i'=s_i$，则 $\Delta E=0$；若 $s_i'=+1$（此时 $s_i=-1$ 且 $h_i>0$），则 $\Delta E=-2h_i<0$；若 $s_i'=-1$（此时 $s_i=+1$ 且 $h_i<0$），则 $\Delta E=2h_i<0$。故 $\Delta E\le0$。

(ii) 状态空间 $\{-1,+1\}^N$ 有限（$2^N$ 个状态），能量函数有下界，每次更新严格下降或不变，故必在有限步内到达能量不再下降的状态——即局部极小点。

(iii) 局部极小点满足 $s_i=\operatorname{sgn}(h_i)$ 对所有 $i$，即不动点条件 $f(\mathbf{s})=\mathbf{s}$。其吸引域由所有收敛到该点的初始状态组成（由 (ii) 的收敛性保证）。
\end{proof}

\begin{corollary}[Hopfield 网络存储容量]\label{cor:hopfield_capacity}
Hopfield 网络可稳定存储的模式数为 $M\approx0.138N$（$N$ 为神经元数）。当存储模式数超过此容量时，模式间的干扰导致吸引子融合或产生虚假吸引子（Spurious Attractors）。
\end{corollary}

\begin{proof}
由统计物理方法，当 $M$ 个随机模式被 Hebb 规则存储时，模式 $i$ 处的信噪比为 $O(\sqrt{N/M})$。当 $M>0.138N$ 时，信噪比下降使模式不再是能量极小点或吸引域被压缩到可忽略。详细推导见 Hertz--Krogh--Palmer, \textit{Introduction to the Theory of Neural Computation}, Addison-Wesley, 1991, Ch. 2。
\end{proof}

\begin{theorem}[循环神经网络的极限环吸引子]\label{thm:hopfield_limit_cycle}
设循环神经网络的状态更新为
\[
\mathbf{h}_{t+1}=\sigma(\mathbf{W}\mathbf{h}_t+\mathbf{U}\mathbf{x}_t+\mathbf{b}),
\]
其中 $\sigma$ 为逐点激活函数，$\mathbf{W}$ 为循环权重。若 $\mathbf{W}$ 的谱半径略大于 $1$ 且激活函数为有界非线性，则网络存在\textbf{极限环吸引子}——即周期吸引子。
\end{theorem}

\begin{proof}
由非线性动力系统的分岔理论：当循环权重谱半径从小于 $1$ 增加到大于 $1$ 时，系统经历 Hopf 分岔——不动点失稳并产生极限环。激活函数有界保证系统状态有界，故极限环为吸引子。周期由 $\mathbf{W}$ 的虚部特征值频率与非线性强度共同决定。
\end{proof}

\begin{proposition}[极限环吸引子的周期条件]\label{prop:limit_cycle_period}
设循环神经网络进入极限环吸引子，其周期为 $T_{\text{lc}}$。则 $T_{\text{lc}}$ 与网络参数的对应关系为
\[
T_{\text{lc}}=\frac{2\pi}{|\operatorname{Im}(\lambda_W)|},
\]
其中 $\lambda_W$ 是 $\mathbf{W}$ 在 Hopf 分岔点处的临界特征值对。
\end{proposition}

\begin{proof}
在 Hopf 分岔附近，系统的线性化动力学由特征值 $\lambda_{1,2}=\alpha\pm i\omega$ 描述，其中 $\alpha$ 为实部（过零时为分岔点），$\omega$ 为虚部。极限环的角频率为 $\omega$，对应周期 $T_{\text{lc}}=2\pi/\omega=2\pi/|\operatorname{Im}(\lambda_W)|$。
\end{proof}

\begin{example}[连续吸引子与记忆保持]\label{ex:continuous_attractor}
在 RNN 中，若循环权重 $\mathbf{W}$ 的特征值全为 $1$（临界情形），则系统存在\textbf{连续吸引子}——吸引子不是孤立点，而是一整条曲线或曲面。连续吸引子可用于保持工作记忆：初始条件沿吸引子流形滑动，不衰减也不发散——这是神经网络实现“持久活动”的动力学机制。
\end{example}

\begin{remark}[吸引子与守恒量的对比]\label{rem:attractor_vs_conservation}
吸引子与守恒量在周期结构中的角色对比如下：
\begin{enumerate}[label=(\roman*)]
    \item \textbf{守恒量}：约束轨道保持在低维流形上（\cref{thm:regular_level_set}），但不直接产生周期；
    \item \textbf{吸引子}：通过耗散把轨道吸引到特定的低维子集，极限环吸引子直接给出周期。
\end{enumerate}
在耗散系统中，吸引子代替守恒量约束轨道（参见「周期基础」章对守恒与周期关系的讨论）。这与「周期系统」章注 \ref{rem:ps_conservation_role} 中“守恒量帮助发现周期但不直接产生周期”的立场一致——真正产生周期的代数枢纽是有限阶作用（离散情形）或吸引子（连续耗散情形）。
\end{remark}

\begin{remark}[吸引子与有限阶作用的关系]\label{rem:attractor_finite_order}
在有限状态 RNN 中，极限环吸引子与「周期基础」章的有限阶作用直接对应：
\begin{enumerate}[label=(\roman*)]
    \item 离散状态空间有限时，每条极限环对应一个循环群 $C_{T_{\text{lc}}}$；
    \item 极限环的周期 $T_{\text{lc}}$ 即有限阶作用的阶数；
    \item 吸引域的划分对应「周期基础」章中不同轨道的分类。
\end{enumerate}
这一对应说明：智能体的动态行为在离散情形下可完全纳入「周期基础」章的有限阶作用框架。
\end{remark}

\section{推理的周期迭代结构}
\label{sec:reasoning_periodicity}

\begin{definition}[推理链与推理映射]\label{def:reasoning_chain}
设智能体的推理状态空间为 $\mathcal{R}\subseteq\mathbb{R}^{d_r}$，推理问题 $q$ 的初始状态为 $\mathbf{s}_0=\psi_{\text{init}}(q)$。\textbf{推理映射}为
\[
\mathcal{T}:\mathcal{R}\times\mathcal{Q}\to\mathcal{R},\qquad
(\mathbf{s},q)\mapsto\mathcal{T}(\mathbf{s},q),
\]
其中 $\mathcal{Q}$ 为问题空间。\textbf{推理链}为迭代序列
\[
\mathbf{s}_{t+1}=\mathcal{T}(\mathbf{s}_t,q),\qquad t=0,1,2,\dots
\]
推理链在 $T_{\text{ans}}$ 步后终止于答案状态 $\mathbf{s}_{T_{\text{ans}}}=\psi_{\text{ans}}^{-1}(\text{答案})$。
\end{definition}

\begin{theorem}[推理的迭代闭合条件]\label{thm:reasoning_fixed_point}
设推理映射 $\mathcal{T}(\cdot,q):\mathcal{R}\to\mathcal{R}$ 为 $L$-压缩映射（$0\le L<1$）。则：
\begin{enumerate}[label=(\roman*)]
    \item 推理链存在唯一不动点 $\mathbf{s}^*=\mathcal{T}(\mathbf{s}^*,q)$——即唯一答案；
    \item 推理链以几何速率收敛：$\|\mathbf{s}_t-\mathbf{s}^*\|\le L^t\|\mathbf{s}_0-\mathbf{s}^*\|$；
    \item 达到精度 $\varepsilon$ 所需推理步数为 $T_{\text{ans}}=O(\log(1/\varepsilon))$。
\end{enumerate}
\end{theorem}

\begin{proof}
由 Banach 不动点定理（\cref{thm:banach_fixed_point_algo}），压缩映射在完备度量空间上有唯一不动点，且迭代以 $L$-几何速率收敛。达到精度 $\varepsilon$ 需 $L^{T_{\text{ans}}}\le\varepsilon/\|\mathbf{s}_0-\mathbf{s}^*\|$，取对数得 $T_{\text{ans}}\ge\log(1/\varepsilon)/\log(1/L)$。
\end{proof}

\begin{proposition}[多步推理的相位累积]\label{prop:reasoning_phase_accumulation}
设推理映射在复数化表示下为相位演化
\[
\theta_{t+1}=\theta_t+\Delta\theta(\mathbf{s}_t,q),
\]
其中 $\Delta\theta(\mathbf{s}_t,q)$ 为第 $t$ 步推理的相位增量。若 $\Delta\theta$ 在收敛过程中趋于常数 $\Delta\theta^*$，则推理链的相位以 $\Delta\theta^*$ 为步长累积，$T_{\text{ans}}$ 步后总相位为
\[
\Phi_{\text{total}}=T_{\text{ans}}\cdot\Delta\theta^*.
\]
当 $\Phi_{\text{total}}\in2\pi\mathbb{Z}$ 时，推理链在相位空间中\textbf{闭合}——即完成一个完整推理周期。
\end{proposition}

\begin{proof}
由相位演化方程迭代：$\theta_t=\theta_0+\sum_{k=0}^{t-1}\Delta\theta(\mathbf{s}_k,q)$。当 $\Delta\theta(\mathbf{s}_k,q)\to\Delta\theta^*$ 时，累积相位 $\Phi_{\text{total}}\approx T_{\text{ans}}\cdot\Delta\theta^*$。闭合条件 $\Phi_{\text{total}}\equiv0\pmod{2\pi}$ 即推理链在相位空间回到起点。
\end{proof}

\begin{lemma}[推理链深度与信息传递]\label{lem:reasoning_depth_information}
设推理链的信息传递效率为 $\rho\in(0,1)$（每步保留的比例）。$T$ 步推理后剩余信息比例为 $\rho^T$。为使最终答案的信息量不低于初始信息量的 $\varepsilon$ 倍，需要
\[
T\le\frac{\log\varepsilon}{\log\rho}.
\]
\end{lemma}

\begin{proof}
$T$ 步后信息比例为 $\rho^T$。要求 $\rho^T\ge\varepsilon$，取对数得 $T\log\rho\ge\log\varepsilon$。因 $\log\rho<0$，除以 $\log\rho$ 反转不等号，得 $T\le\log\varepsilon/\log\rho$。
\end{proof}

\begin{corollary}[推理链的最优深度]\label{cor:optimal_reasoning_depth}
推理链的最优深度是\textbf{信息传递}与\textbf{计算成本}的权衡：
\begin{enumerate}[label=(\roman*)]
    \item 深度太小：信息丢失过多，答案精度低；
    \item 深度太大：信息高度保留但计算成本高，且可能因数值误差累积而恶化。
\end{enumerate}
最优深度为 $T^*=O(\log(1/\varepsilon)/\log(1/\rho))$——与压缩推理映射的收敛步数同阶。
\end{corollary}

\begin{proof}
综合 \cref{thm:reasoning_fixed_point} 与 \cref{lem:reasoning_depth_information}：两者都给出 $T=O(\log(1/\varepsilon))$ 的步数尺度，前者来自压缩映射的几何收敛，后者来自信息传递的指数衰减。
\end{proof}

\begin{theorem}[Chain-of-Thought 推理的相位锁定]\label{thm:cot_phase_locking}
设 Chain-of-Thought（CoT）推理包含 $T_{\text{CoT}}$ 个中间步骤。若每个中间步骤的相位更新为 $\Delta\theta$（常数），则 CoT 推理链的相位锁定条件为
\[
T_{\text{CoT}}\cdot\Delta\theta\in2\pi\mathbb{Z}.
\]
此时推理链在相位空间中闭合，对应一个“推理周期”。
\end{theorem}

\begin{proof}
CoT 推理链的相位累积为 $\Phi_{\text{total}}=T_{\text{CoT}}\cdot\Delta\theta$（由 \cref{prop:reasoning_phase_accumulation}）。闭合条件 $\Phi_{\text{total}}\in2\pi\mathbb{Z}$ 即得。
\end{proof}

\begin{example}[数学证明的相位结构]\label{ex:math_proof_phase}
在数学证明推理中，每一步推理对应一个逻辑蕴含。若每一步的逻辑强度（相位增量）恒定，则证明完成时逻辑链的相位应闭合（即从假设推出结论，总相位累积为 $2\pi$）。这是“推理闭合”的相位解释：\textbf{证明完成等价于推理链在逻辑相位空间中的闭合}。
\end{example}

\begin{remark}[推理与不动点迭代的对应]\label{rem:reasoning_fixed_point_correspondence}
推理与算法系统章的不动点迭代（\cref{def:fixed_point_iteration}）在数学结构上完全对应：
\begin{enumerate}[label=(\roman*)]
    \item 推理映射 $\mathcal{T}$ 对应迭代映射 $\Phi$；
    \item 推理答案 $\mathbf{s}^*$ 对应不动点 $\mathbf{x}^*$；
    \item 推理收敛条件对应压缩性条件 $L<1$；
    \item CoT 的深度 $T_{\text{CoT}}$ 对应迭代步数 $T$。
\end{enumerate}
因此，「算法系统之周期结构」章的收敛分析可直接应用于推理链的收敛性分析。
\end{remark}

\section{知识表示的周期分析}
\label{sec:knowledge_periodicity}

\begin{definition}[知识图谱与关系嵌入]\label{def:knowledge_graph_embedding}
\textbf{知识图谱}是三元组集合 $\mathcal{G}=\{(h,r,t)\}$，其中 $h$ 为头实体，$t$ 为尾实体，$r$ 为关系。\textbf{关系嵌入}为映射
\[
\varphi_{\text{KG}}:\mathcal{E}\cup\mathcal{R}\to\mathbb{R}^{d_k},
\]
其中 $\mathcal{E}$ 为实体集，$\mathcal{R}$ 为关系集。经典 TransE 模型的嵌入目标为
\[
\mathbf{e}_h+\mathbf{e}_r\approx\mathbf{e}_t.
\]
\end{definition}

\begin{theorem}[关系嵌入的平移-旋转分解]\label{thm:relation_translation_rotation}
设知识图谱嵌入满足 TransE 关系 $\mathbf{e}_h+\mathbf{e}_r\approx\mathbf{e}_t$。在复数化表示下，该关系等价于
\[
\mathbf{e}_t=\mathbf{e}_h\cdot e^{i\Delta\theta_r}\quad\text{或}\quad
\mathbf{e}_t=\mathbf{e}_h+\mathbf{u}_r,
\]
其中 $e^{i\Delta\theta_r}$ 为关系的相位旋转，$\mathbf{u}_r$ 为关系的平移向量。二者分别对应\textbf{旋转表示}（RotatE）与\textbf{平移表示}（TransE）。
\end{theorem}

\begin{proof}
（平移）由 TransE 目标直接读得 $\mathbf{e}_t-\mathbf{e}_h\approx\mathbf{e}_r$——关系为平移向量。
（旋转）将嵌入复数化 $e^{i\theta_h}$，取关系相位 $\Delta\theta_r$，则尾实体相位 $\theta_t=\theta_h+\Delta\theta_r$，即 $\mathbf{e}_t=\mathbf{e}_h\cdot e^{i\Delta\theta_r}$。两种表示在模长近似相等且 $\Delta\theta_r$ 小时对偶。
\end{proof}

\begin{proposition}[知识推理的路径周期]\label{prop:knowledge_path_period}
设知识图谱中存在一条关系路径
\[
h\xrightarrow{r_1}e_1\xrightarrow{r_2}e_2\xrightarrow{r_3}\cdots\xrightarrow{r_n}t.
\]
则路径的总关系相位为
\[
\Delta\Phi_{\text{path}}=\sum_{k=1}^n\Delta\theta_{r_k}.
\]
当 $\Delta\Phi_{\text{path}}=0\pmod{2\pi}$ 时，路径闭合——即 $h$ 与 $t$ 是同一实体的不同关系路径。
\end{proposition}

\begin{proof}
由 \cref{thm:relation_translation_rotation} 的旋转表示，每一步关系的相位累积为 $\theta_{e_{k}}=\theta_{e_{k-1}}+\Delta\theta_{r_k}$。迭代得 $\theta_t=\theta_h+\sum_k\Delta\theta_{r_k}$。故路径闭合 $\theta_t=\theta_h$ 等价于 $\sum_k\Delta\theta_{r_k}\equiv0\pmod{2\pi}$。
\end{proof}

\begin{lemma}[传递关系的相位合成]\label{lem:transitive_relation_phase}
设关系 $r_1$ 与 $r_2$ 的相位增量分别为 $\Delta\theta_{r_1}$ 与 $\Delta\theta_{r_2}$。若 $r_3=r_1\circ r_2$（$r_3$ 为 $r_1$ 与 $r_2$ 的复合关系），则
\[
\Delta\theta_{r_3}=\Delta\theta_{r_1}+\Delta\theta_{r_2}.
\]
即关系的相位在复合下满足加法——这是「标量与矢量」章复数乘法的相位相加（\cref{prop:complex_multiplication}）在知识推理中的实现。
\end{lemma}

\begin{proof}
由 \cref{prop:knowledge_path_period} 中两步路径的推导：$\theta_{e_2}=\theta_h+\Delta\theta_{r_1}+\Delta\theta_{r_2}$。复合关系 $r_3$ 的相位增量定义为 $\theta_{e_2}-\theta_h=\Delta\theta_{r_1}+\Delta\theta_{r_2}$。
\end{proof}

\begin{example}[TransE 嵌入的平移结构]\label{ex:transe_translation}
在 TransE 模型中，关系“首都”的嵌入为平移向量 $\mathbf{e}_{\text{首都}}$，使
\[
\mathbf{e}_{\text{中国}}+\mathbf{e}_{\text{首都}}\approx\mathbf{e}_{\text{北京}},\quad
\mathbf{e}_{\text{法国}}+\mathbf{e}_{\text{首都}}\approx\mathbf{e}_{\text{巴黎}}.
\]
这体现了\textbf{关系的平移不变性}：同一关系在不同实体对之间表现为\textbf{相同的平移向量}——类比关系的平移版本（\cref{prop:embedding_rotation_analogy}）。
\end{example}

\section{学习动力学的吸引子分析}
\label{sec:learning_attractor_analysis}

\begin{definition}[训练的吸引子结构]\label{def:training_attractor}
设参数空间的训练动力学为 $\boldsymbol{\theta}_{t+1}=\boldsymbol{\theta}_t-\eta\nabla\mathcal{L}(\boldsymbol{\theta}_t)$。\textbf{训练的吸引子}为参数空间中的不动点 $\boldsymbol{\theta}^*$（$\nabla\mathcal{L}(\boldsymbol{\theta}^*)=\mathbf{0}$，且 Hessian 正定）；在二次损失下 SGD 只有不动点吸引子——学习率较大时参数带符号交替地几何收敛到 $\boldsymbol{\theta}^*$，这种振荡\textbf{不是}极限环（见 \cref{thm:sgd_limit_cycle}）。
\end{definition}

\begin{theorem}[SGD 的收敛与欠阻尼振荡]\label{thm:sgd_limit_cycle}
设损失函数 $\mathcal{L}(\boldsymbol{\theta})$ 的 Hessian 特征值为 $\{\lambda_i\}$，学习率为 $\eta$。则：
\begin{enumerate}[label=(\roman*)]
    \item 若 $\eta\lambda_i<1$ 对所有 $i$（过阻尼区），则 SGD 收敛到不动点；
    \item 若 $1<\eta\lambda_i<2$ 对某 $i$（欠阻尼区），则 SGD 在该特征方向上带符号交替地振荡，但误差仍几何收敛到 $0$（因 $|1-\eta\lambda_i|<1$），参数最终收敛到不动点 $\boldsymbol{\theta}^*$；
    \item 若 $1<\eta\lambda_i<2$ 对所有 $i$，则各特征方向的误差均带符号交替地几何收敛到 $0$，参数以衰减振荡的方式收敛到 $\boldsymbol{\theta}^*$。二次损失下\textbf{不存在极限环}：\textbf{振荡不等于极限环}——极限环要求轨道闭合且孤立，而此处 $|1-\eta\lambda_i|<1$ 使振幅严格衰减。
\end{enumerate}
\end{theorem}

\begin{proof}
由 \cref{ch32_gd_eigen}，第 $i$ 个特征方向的误差演化为 $e_t^{(i)}=(1-\eta\lambda_i)^te_0^{(i)}$。若 $0<\eta\lambda_i<1$，则 $0<1-\eta\lambda_i<1$——误差单调衰减到零。若 $1<\eta\lambda_i<2$，则 $-1<1-\eta\lambda_i<0$——误差符号每步交替而模严格减小，故参数在极小值附近\textbf{带符号交替地}衰减到 $\boldsymbol{\theta}^*$。多特征方向同时欠阻尼时，参数轨迹是绕极小值的衰减振荡（各方向交替变号），其振幅按 $\prod_i|1-\eta\lambda_i|$ 逐类衰减，\textbf{并不形成极限环}：极限环要求振幅不衰减且轨道孤立。
\end{proof}

\begin{proposition}[泛化与吸引子宽度的关系]\label{prop:generalisation_attractor_width}
设训练收敛到吸引子 $\boldsymbol{\theta}^*$，其 Hessian 特征值为 $\{\lambda_i\}$。定义吸引子宽度为
\[
W=\frac{1}{d}\sum_{i=1}^{d}\frac{1}{\lambda_i}.
\]
则：
\begin{enumerate}[label=(\roman*)]
    \item 宽吸引子（$W$ 大，$\lambda_i$ 小）对应平坦极小值，泛化性能好；
    \item 窄吸引子（$W$ 小，$\lambda_i$ 大）对应尖锐极小值，泛化性能差。
\end{enumerate}
这是深度学习中的\textbf{平坦极小值假说}。
\end{proposition}

\begin{proof}
平坦极小值附近损失函数变化缓慢，参数在扰动下仍保持低损失；尖锐极小值附近损失函数变化剧烈，小扰动导致损失大幅上升。宽吸引子的 Hessian 特征值小（$W$ 大），对应前者；窄吸引子的特征值大（$W$ 小），对应后者。泛化性能由此依赖吸引子宽度。
\end{proof}

\begin{example}[模式记忆与容量分析]\label{ex:pattern_memory}
Hopfield 网络（\cref{thm:hopfield_convergence}）的每个局部极小点对应一个\textbf{记忆模式}。存储 $M$ 个模式时，吸引子总数为 $M$（加上可能产生的虚假吸引子）。由容量分析（\cref{cor:hopfield_capacity}），当 $M\le0.138N$ 时，真实模式为稳定吸引子；当 $M$ 超过此容量时，模式间干扰使吸引子融合或产生虚假吸引子。这是智能体长期记忆容量的数学上界。
\end{example}

\begin{remark}[SGD 的欠阻尼振荡与泛化]\label{rem:sgd_limit_cycle_generalisation}
SGD 在欠阻尼区带符号交替地收敛到不动点，这一衰减振荡对泛化有重要影响：
\begin{enumerate}[label=(\roman*)]
    \item 欠阻尼振荡中的参数波动相当于\textbf{隐式正则化}——收敛过程不会精确停留在尖锐极小值；
    \item 衰减振荡的振幅衰减率（由 $|1-\eta\lambda_{\max}|$ 决定）与有效学习率 $\eta\lambda_{\max}$ 相关，可作为泛化误差的经验指标；
    \item 这与「算法系统之周期结构」章中“梯度下降的阻尼振荡”共享同一数学结构——训练后期的振荡并非缺陷，而是隐式泛化的来源。
\end{enumerate}
\end{remark}

\section{注意力机制的周期结构}

\begin{definition}[自注意力的周期解释]\label{def:self_attention_periodic}
Transformer 模型的核心机制是\textbf{缩放点积注意力}（Scaled Dot-Product Attention）：
\[
\operatorname{Attention}(\mathbf{Q}, \mathbf{K}, \mathbf{V}) = \operatorname{softmax}\left(\frac{\mathbf{Q}\mathbf{K}^{\mathsf T}}{\sqrt{d_k}}\right) \mathbf{V}.
\]
注意力机制是三种操作在序列建模中的统一实现：
\begin{enumerate}[label=(\roman*)]
    \item \textbf{平移操作}：注意力窗口在输入序列上的滑动是平移操作——Query 向量 $\mathbf{Q}$ 在序列上移动，与不同位置的 Key 向量 $\mathbf{K}$ 计算相似度。注意力的平移类似于「离散采样逼近」章的窗函数在时间序列上的滑动（\cref{def:window_function}）。
    \item \textbf{缩放操作}：多头注意力（Multi-Head Attention）将注意力并行投射到 $h$ 个不同的子空间——每个头关注不同尺度的依赖关系。这是标度分解：不同头对应不同的缩放因子（不同的 $d_k$ 决定了注意力的“分辨率”）。
    \item \textbf{旋转操作}：位置编码（Positional Encoding）使用正弦和余弦函数注入位置信息：$PE_{(pos, 2i)} = \sin(pos/10000^{2i/d})$，$PE_{(pos, 2i+1)} = \cos(pos/10000^{2i/d})$。位置编码是旋转操作在序列建模中的实现——不同频率的正弦波对应嵌套周转圆的不同层级，位置 $pos$ 是相位索引。
\end{enumerate}
\end{definition}

\begin{proposition}[位置编码的傅里叶结构]\label{prop:positional_encoding_fourier}
正弦位置编码可重写为复指数形式：$PE_{pos} = e^{i \omega \cdot pos}$，其中 $\omega_i = 1/10000^{2i/d}$ 是各级频率。不同维度 $i$ 对应不同的角速度 $\omega_i$——位置编码是嵌套周转系统在序列位置上的傅里叶展开。第 $i$ 维位置编码对应半径为 $1$、角速度为 $\omega_i$ 的周转圆（\cref{def:epicycle}），位置 $pos$ 对应旋转角度。

进一步，当注意力权重仅由位置编码贡献时（例如输入特征 $\mathbf{x} = \mathbf{0}$），点积 $\mathbf{q}_i \cdot \mathbf{k}_j$ 给出
\[
\mathbf{q}_i \cdot \mathbf{k}_j = \sum_{l=1}^{d/2} \cos(\omega_l (pos_i - pos_j)),
\]
即\textbf{多频率差角余弦相似度之和}——这正是「标量与矢量」章的基础差角余弦相似度 $\operatorname{cosine}(\mathbf{a},\mathbf{b}) = \cos(\alpha - \beta)$（\cref{def:basic_diff_similarity}）在多频率上的叠加。
\end{proposition}

\begin{proof}
由欧拉公式，$(\sin(pos \cdot \omega_i), \cos(pos \cdot \omega_i))$ 恰为 $e^{i\omega_i \cdot pos}$ 的虚部与实部——位置编码的每一维对是单位圆上的点，角速度 $\omega_i = 1/10000^{2i/d}$。

对点积的显式计算：取 $\mathbf{W}_Q = \mathbf{W}_K = \mathbf{I}$（简化情形），则
\[
\mathbf{q}_i \cdot \mathbf{k}_j = PE_{pos_i} \cdot PE_{pos_j} = \sum_{l} \left[\sin(\omega_l pos_i)\sin(\omega_l pos_j) + \cos(\omega_l pos_i)\cos(\omega_l pos_j)\right].
\]
由积化和差公式，$\sin A\sin B + \cos A\cos B = \cos(A - B)$，故
\[
\mathbf{q}_i \cdot \mathbf{k}_j = \sum_{l} \cos(\omega_l (pos_i - pos_j)).
\]
这正是差角余弦相似度 $\cos(\alpha - \beta)$（取 $\alpha = \omega_l pos_i$，$\beta = \omega_l pos_j$）在多频率上的求和。
\end{proof}

\begin{remark}[位置编码点积的平移不变性]\label{rem:positional_encoding_translation_invariance}
由 \cref{prop:positional_encoding_fourier} 的推导，位置编码的点积相似度只依赖于\textbf{相对位置} $pos_i - pos_j$，与绝对位置无关。这一平移不变性使得位置编码具有“相对位置感知”能力——这与「级数与卷积」章的相位平移（\cref{prop:phase_translation}）在数学结构上一致。
\end{remark}

\section{循环神经网络的周期记忆}

\begin{definition}[LSTM 的门控周期]\label{def:lstm_gates}
长短期记忆网络（LSTM）通过三个门控机制控制信息的流动：
\begin{enumerate}[label=(\roman*)]
    \item \textbf{遗忘门} $f_t = \sigma(\mathbf{W}_f \cdot [\mathbf{h}_{t-1}, \mathbf{x}_t] + \mathbf{b}_f)$：决定哪些旧信息需要被丢弃。$f_t \to 1$ 保持信息，$f_t \to 0$ 遗忘信息。
    \item \textbf{输入门} $i_t = \sigma(\mathbf{W}_i \cdot [\mathbf{h}_{t-1}, \mathbf{x}_t] + \mathbf{b}_i)$：决定哪些新信息需要被写入。
    \item \textbf{输出门} $o_t = \sigma(\mathbf{W}_o \cdot [\mathbf{h}_{t-1}, \mathbf{x}_t] + \mathbf{b}_o)$：决定哪些信息需要被输出到隐状态。
\end{enumerate}
LSTM 的门控机制是旋转操作在隐状态空间中的实现：遗忘门 $f_t$ 控制旧信息的衰减速率，输入门 $i_t$ 控制新信息的写入速率，输出门 $o_t$ 控制信息从细胞状态到隐状态的投影。三个门的周期性开闭驱动了隐状态 $\mathbf{h}_t$ 在状态空间中的旋转。
\end{definition}

\begin{proposition}[LSTM 的长程周期保持]\label{prop:lstm_long_range}
LSTM 能够学习序列中的长程周期依赖。设输入序列包含周期为 $T$ 的周期模式。LSTM 的遗忘门 $f_t$ 在训练中学习到以 $T$ 为周期的开闭模式——在周期的关键位置 $f_t \approx 1$（保持信息），在周期的非关键位置 $f_t \approx 0$（遗忘信息）。细胞状态 $\mathbf{c}_t$ 的更新 $\mathbf{c}_t = f_t \odot \mathbf{c}_{t-1} + i_t \odot \tilde{\mathbf{c}}_t$ 等价于在状态空间中进行选择性记忆——$f_t$ 的周期开闭类似于采样门，以周期 $T$ 对输入进行采样。

LSTM 的长程记忆是相位锁定的神经网络实现：遗忘门的周期性开闭与输入序列的周期同步——外部驱动（输入序列）与内部周期（LSTM 隐状态）的\textbf{相位锁定}（\cref{prop:phase-locking}）。当遗忘门的开闭周期与输入序列的周期匹配时，LSTM 实现最优记忆保持。
\end{proposition}

\begin{proof}
考虑简化的线性情形：遗忘门 $f_t = f$ 为常数、无输入门调制（$i_t \cdot \tilde{\mathbf{c}}_t = \mathbf{u}_t$）。细胞状态递推 $\mathbf{c}_t = f \mathbf{c}_{t-1} + \mathbf{u}_t$ 的解为
\[
\mathbf{c}_t = \sum_{k=0}^{t} f^{t-k} \mathbf{u}_k,
\]
即以 $f$ 为底数的指数加权和，记忆有效长度 $\tau_{\text{mem}} = -1/\ln f$（$f \to 1$ 时记忆长度发散）。

当遗忘门周期开闭 $f_t = f(t \bmod T)$（周期 $T$）时，乘积 $\prod_{k=0}^{T-1} f_k$ 在 $T$ 步内回到初值——周期记忆的\textbf{临界保持条件}为 $\prod_{k=0}^{T-1} f_k = 1$。此时，周期分量（频率与 $1/T$ 公度者）无衰减传播；非周期分量被逐周期压缩。

这正是相位锁定（\cref{prop:phase-locking}）：遗忘门的相位（由 $f_t$ 的周期结构决定）与序列的相位（由输入周期决定）一致时记忆最优。
\end{proof}

\section{强化学习的周期探索}

\begin{definition}[探索-利用的周期交替]\label{def:exploration_exploitation}
强化学习中智能体面临探索（尝试新动作以获取更多信息）和利用（执行已知最优动作以获取最大回报）的经典权衡。$\varepsilon$-贪婪策略以概率 $\varepsilon$ 随机探索，以概率 $1-\varepsilon$ 贪婪利用。$\varepsilon$ 通常随时间衰减：
\[
\varepsilon(t) = \varepsilon_{\min} + (\varepsilon_{\max} - \varepsilon_{\min}) e^{-t/\tau},
\]
其中 $\tau$ 是衰减时间常数。$\varepsilon$ 的指数衰减是标度衰减在强化学习中的实现——探索强度随时间按标度因子 $e^{-1/\tau}$ 衰减。在训练的早期（$t \ll \tau$），$\varepsilon \approx \varepsilon_{\max}$——智能体大量探索（随机旋转）。在训练的后期（$t \gg \tau$），$\varepsilon \approx \varepsilon_{\min}$——智能体主要利用（确定性平移）。

探索过程中的随机源由 PRNG 提供（\cref{def:prng_period}）——其周期长度决定了探索随机性的质量。若 PRNG 周期过短，探索策略会在短时间内重复，降低探索效率。
\end{definition}

\begin{proposition}[策略迭代的周期收敛]\label{prop:policy_iteration_cycle}
策略迭代（Policy Iteration）在策略评估和策略改进之间交替：
\[
\pi_0 \xrightarrow{\text{评估}} Q^{\pi_0} \xrightarrow{\text{改进}} \pi_1 \xrightarrow{\text{评估}} Q^{\pi_1} \xrightarrow{\text{改进}} \pi_2 \to \cdots \to \pi^*.
\]
每次策略改进使策略向最优策略 $\pi^*$ 旋转一个角度——策略在策略空间中绕次优策略旋转，逐步逼近最优策略。策略评估对应沿值函数方向移动，策略改进对应沿贪婪方向旋转策略。策略迭代的收敛是梯度系统向平衡点的收敛——最优策略 $\pi^*$ 是策略空间中的吸引子。
\end{proposition}

\begin{proof}
（有限 MDP 情形）\textbf{策略改进定理}：若 $\pi'$ 为关于 $Q^{\pi}$ 的贪婪策略（$\pi'(s) = \arg\max_a Q^{\pi}(s,a)$），则 $V^{\pi'}(s) \ge V^{\pi}(s)$ 对所有 $s$ 成立。证明：由 $V^{\pi}(s) \le Q^{\pi}(s, \pi'(s))$ 与贝尔曼期望递推逐状态验证——对任意策略 $\pi$，
\[
\begin{aligned}
V^{\pi}(s) &= \mathbb{E}_{a \sim \pi(\cdot|s)}\left[ r(s,a) + \gamma \sum_{s'} P(s'|s,a) V^{\pi}(s') \right]\\
&\le \max_a \left[ r(s,a) + \gamma \sum_{s'} P(s'|s,a) V^{\pi}(s') \right] = Q^{\pi}(s, \pi'(s)).
\end{aligned}
\]
反复应用贝尔曼算子，$\pi'$ 的值函数逐点不低于 $\pi$。

有限状态-动作空间中策略有限，故单调改进序列 $\pi_0 \preceq \pi_1 \preceq \dots$ 必在有限步达到不动点 $\pi^*$（最优贝尔曼方程的解）——策略迭代在策略空间中的“旋转”最终闭合于吸引子 $\pi^*$。
\end{proof}

\section{多智能体系统的相位同步}

\begin{definition}[多智能体强化学习的策略同步]\label{def:multi_agent_sync}
在多智能体强化学习（MARL）中，$N$ 个智能体各自维护一个策略 $\pi_i$。智能体之间通过共享参数、共享经验或通信实现策略协调。当智能体的奖励函数具有协同性（所有智能体的奖励正相关）时，策略参数在参数空间中趋于\textbf{相位同步}——各智能体的策略参数向量 $\boldsymbol{\theta}_i$ 趋于一致方向。

多智能体策略同步是相位锁定的多体实现：每个智能体是一个独立的周期子系统（策略更新的周期振荡），智能体之间的通信是子结构之间的耦合系数 $c_{ij}$。当耦合足够强时，所有智能体的策略参数趋于相同的相位——这与 Kuramoto 模型中振子的相位同步（\cref{prop:kuramoto_model}）在数学结构上一致，其同步阈值由 Kuramoto 临界耦合强度 $K_c = 2/(\pi g(\Omega))$ 给出（$g$ 为固有频率分布，$\Omega$ 为其众数）。
\end{definition}

\begin{proposition}[参数平均的周期收敛]\label{prop:parameter_averaging}
在分布式训练中，多个工作节点（worker）独立计算梯度，参数服务器（parameter server）周期性地平均各节点的参数：
\[
\boldsymbol{\theta}^{(t+1)} = \frac{1}{N} \sum_{i=1}^{N} \boldsymbol{\theta}_i^{(t)}.
\]
参数平均的周期 $T_{\text{sync}}$（每隔多少步同步一次）是同步频率参数。$T_{\text{sync}} = 1$（每步同步）对应强耦合——训练稳定但通信开销大。$T_{\text{sync}} \gg 1$（多步异步）对应弱耦合——通信开销小但各节点参数可能发散（相位漂移）。最优同步周期需要在收敛速度和通信效率之间权衡。
\end{proposition}

\begin{proof}
设各节点独立执行 $T_{\text{sync}}$ 步梯度下降后平均参数。单节点 $T_{\text{sync}}$ 步的误差由 \cref{ch32_gd_eigen} 按 $(1-\eta\lambda_i)^{T_{\text{sync}}}$ 收缩；平均操作是凸组合，不增大各节点与全局均值的分散度（凸性）。故 $T_{\text{sync}}$ 步内收敛速度约为同步执行（$T_{\text{sync}}=1$）的 $\frac{1}{T_{\text{sync}}}\ln\frac{1}{|1-\eta\lambda_i|}$ 倍，而通信量降为 $1/T_{\text{sync}}$——同步周期是在收敛速率与通信成本之间的权衡参数。

这一权衡结构与 Kuramoto 模型的耦合强度权衡（\cref{prop:kuramoto_model}）在数学结构上一致：耦合越强（$T_{\text{sync}}$ 越小）同步越快但通信开销越大；耦合越弱（$T_{\text{sync}}$ 越大）通信越经济但收敛越慢，且当耦合低于 Kuramoto 阈值时同步失效。
\end{proof}

\section{智能体系统的约束关系}

\begin{theorem}[智能体系统的三种约束关系]\label{thm:agent_conservation}
智能体系统的动力学由三种约束关系支配：
\begin{enumerate}[label=(\roman*)]
    \item \textbf{参数约束}：在无正则化的训练中，模型参数的总量（$L^2$ 范数）在优化中不直接受约束。但权重衰减（weight decay）引入参数向原点的梯度拉力——$\boldsymbol{\theta}_{t+1} = \boldsymbol{\theta}_t - \eta(\nabla\mathcal{L} + \lambda \boldsymbol{\theta}_t)$。权重衰减是向平衡点的回归——类似于弹簧的恢复力。
    \item \textbf{梯度传播约束}：在批归一化（Batch Normalisation）中，每层输入的均值和方差在训练中保持——$\mathbb{E}[\mathbf{x}] = 0, \operatorname{Var}[\mathbf{x}] = 1$。梯度在各层之间的传播满足能量约束——梯度的总能量（$L^2$ 范数）在反向传播中不衰减不爆炸。残差连接（Residual Connection）通过恒等映射 $\mathbf{x}_{l+1} = \mathbf{x}_l + \mathcal{F}(\mathbf{x}_l)$ 保证了梯度传播——$\partial \mathbf{x}_{l+1}/\partial \mathbf{x}_l = \mathbf{I} + \partial \mathcal{F}/\partial \mathbf{x}_l$，单位矩阵保证了梯度的无损传播。
    \item \textbf{注意力归一化约束}：注意力权重的总和为 $1$——$\sum_j \alpha_{ij} = 1$ 对所有 $i$ 成立。注意力的归一化是相位锁定的实现——智能体在任意时刻的总注意力是固定的，注意力在不同输入位置之间的分配是相位的旋转。总注意资源在多个注意力头之间的分配构成了注意力归一化约束。
\end{enumerate}
三种约束关系共同保证了智能体系统的训练稳定性和泛化能力——参数约束防止过拟合，梯度传播约束防止梯度消失/爆炸，注意力归一化约束保证信息聚焦的合理性。
\end{theorem}

\begin{proof}
(i) \textbf{参数约束}：权重衰减的连续时间形式为 $d\boldsymbol{\theta}/dt = -\nabla\mathcal{L} - \lambda \boldsymbol{\theta}$。在平衡点 $\nabla\mathcal{L} + \lambda \boldsymbol{\theta}^* = \mathbf{0}$ 附近，$\|\boldsymbol{\theta}^*\| = \|\nabla\mathcal{L}\|/\lambda$——参数的 $L^2$ 范数由损失梯度和衰减系数共同决定。权重衰减引入的是向平衡点的恢复力——类似于弹簧振子中的 $-kx$ 项（\cref{def:spring_oscillator}）。

(ii) \textbf{梯度传播约束}：残差网络的第 $l$ 层输出为 $\mathbf{x}_{l+1} = \mathbf{x}_l + \mathcal{F}(\mathbf{x}_l, \mathbf{W}_l)$。反向传播的梯度为
\[
\frac{\partial \mathcal{L}}{\partial \mathbf{x}_l} = \frac{\partial \mathcal{L}}{\partial \mathbf{x}_{l+1}} \left(\mathbf{I} + \frac{\partial \mathcal{F}}{\partial \mathbf{x}_l}\right).
\]
单位矩阵 $\mathbf{I}$ 保证了梯度可以直接从深层传播到浅层——即使在 $\frac{\partial \mathcal{F}}{\partial \mathbf{x}_l}$ 很小的情况下，梯度也不会消失。这是梯度传播约束在深度网络中的实现。

(iii) \textbf{注意力归一化约束}：$\operatorname{softmax}$ 函数保证了 $\sum_j \alpha_{ij} = 1$——总注意力在所有位置上的分配是归一化的。当某个位置的注意力增加时，其他位置的注意力必然减少。多头部注意力将总注意力分配到 $h$ 个头中，每个头独立进行归一化——总注意力在头之间的分配也是归一化的。这一约束与「概率系统之周期结构」章的概率归一化（\cref{prop:probability_measure_properties}）共享同一个“总量守恒”的数学结构。
\end{proof}

\section{开放问题与研究纲领}
\label{sec:agent_open_problems}

\begin{remark}[开放问题]\label{rem:agent_open_problems}
以下问题超出本章已建立的定理范围，列为\textbf{开放问题}：
\begin{enumerate}[label=(\roman*)]
    \item \textbf{KDI 在非凸损失下的全局收敛性}：\cref{prop:kdi_eaf_convergence} 仅给出定性描述；KDI 的全局收敛速率、逃离鞍点的期望时间等尚缺乏严格结果。
    \item \textbf{CEAF 在深度网络中的表达力}：\cref{thm:universal-approx-epicycle} 是单隐层结果；多层 CEAF 网络是否比单层具有本质优势，是否比 ReLU 网络具有更高的参数效率，尚缺乏严格分离下界。
    \item \textbf{位置编码与周期卷积的精确对应}：\cref{prop:positional_encoding_fourier} 给出了位置编码点积的多频率余弦分解；这一分解是否能完整对应「级数与卷积」章的周期卷积框架（\cref{def:periodic_convolution}），特别是权重核 $W$ 与频率合成算子 $\Omega_{n,m}$ 的对应，尚需进一步研究。
    \item \textbf{多智能体同步与 Kuramoto 模型的精确对应}：\cref{prop:parameter_averaging} 给出了参数平均的定性收敛；能否给出与 Kuramoto 阈值 $K_c$（\cref{prop:kuramoto_model}）精确对应的同步阈值，尚需深入。
    \item \textbf{强化学习探索策略的最优衰减率}：$\varepsilon(t) = \varepsilon_{\min} + (\varepsilon_{\max} - \varepsilon_{\min}) e^{-t/\tau}$ 中的时间常数 $\tau$ 的最优选取，作为 MDP 参数的函数，尚无普遍公式。
    \item \textbf{词嵌入相位结构的定量刻画}：\cref{prop:embedding_rotation_analogy} 给出了类比关系的相位偏移描述；这一偏移是否在所有周期性语义维度上都可观察、是否受训练数据分布的定量控制，尚需实证研究。
    \item \textbf{推理链最优深度的普遍公式}：\cref{cor:optimal_reasoning_depth} 给出了最优深度的数量级，但具体的比例常数与问题的内在复杂度之间的关系尚无严格定理。
    \item \textbf{吸引子宽度与泛化误差的定量关系}：\cref{prop:generalisation_attractor_width} 给出了宽度-泛化的定性关系；定量泛化误差界与吸引子宽度的精确联系尚待建立。
\end{enumerate}
这些问题的解决将把智能体周期结构理论从“有定性保证的框架”提升为“有定量保证的学习算法”。
\end{remark}

\begin{remark}[与全书的统一]\label{rem:agent_chapter_unity}
本章建立的智能体系统周期结构理论是周期数学在人工智能中的系统应用，同时是“数据分析 $\to$ 激活函数 $\to$ 算法 $\to$ 智能体”四章链条的汇聚点。

\textbf{来自前三章的继承}：
\begin{enumerate}[label=(\roman*)]
    \item 梯度下降的振荡收敛（\cref{ch32_gd_eigen}）是平移操作在参数空间中的实现——学习率决定了平移步长，Hessian 特征值决定了各方向的收敛速率；
    \item 动量法（\cref{prop:momentum_rotation}）在参数空间中引入旋转分量——参数轨迹绕损失地形的沟壑旋转加速收敛；
    \item 学习率调度（\cref{def:lr_scheduler}）是标度衰减在优化中的实现；
    \item 牛顿法（\cref{prop:newton_quadratic}）的雅可比矩阵提供了旋转修正；
    \item 词嵌入的相位表示（\cref{lem:embedding_phase_representation}）与旋转类比关系（\cref{prop:embedding_rotation_analogy}）继承「标量与矢量」章的复指数差角算子（\cref{def:complex_diff_sum}）。
\end{enumerate}

\textbf{本章的贡献}：
\begin{enumerate}[label=(\roman*)]
    \item 注意力机制（\cref{def:self_attention_periodic}）是三种操作在序列建模中的统一——Query-Key 的点积是相位相似度计算（旋转操作），多头并行的多尺度是标度分解（缩放操作），窗口滑动是位置平移（平移操作）；
    \item LSTM 的门控周期（\cref{prop:lstm_long_range}）是旋转操作在隐状态空间中的实现——遗忘门和输入门的周期开闭驱动隐状态的旋转更新；
    \item 强化学习的探索-利用交替（\cref{def:exploration_exploitation}）是缩放操作和旋转操作的耦合；
    \item 多智能体的策略同步（\cref{prop:parameter_averaging}）是相位锁定（\cref{prop:phase-locking}）与 Kuramoto 模型（\cref{prop:kuramoto_model}）的工程实现；
    \item 智能体系统的三种约束关系（\cref{thm:agent_conservation}）——参数约束、梯度传播约束、注意力归一化约束——共同构成了深度学习中训练稳定性的数学保证；
    \item 词嵌入的相位结构（\cref{sec:word_embedding_periodicity}）为自然语言处理中的语义类比提供了周期数学框架；
    \item 吸引子动力学（\cref{sec:attractor_dynamics}）把「周期基础」章的有限阶作用与连续耗散系统的极限环统一到智能体的记忆与振荡结构；
    \item 推理的周期迭代结构（\cref{sec:reasoning_periodicity}）把「算法系统之周期结构」章的不动点迭代推广到推理链的收敛性与闭合性分析；
    \item 知识表示的周期分析（\cref{sec:knowledge_periodicity}）把关系嵌入的平移-旋转对偶纳入周期数学框架。
\end{enumerate}

\textbf{四章链条的完整闭合}：从数据分析的参数提取（\cref{prop:data_to_activation_map}），到激活函数的可微组件（\cref{prop:activation_to_algorithm_transfer}），到算法收敛结构的继承（\cref{prop:algorithm_to_agent_inheritance}），再到智能体训练动力学的完整实现——四章以“嵌套周转系统”为统一数学对象，以“平移-缩放-旋转”三种基本操作为统一语言，完成了从数据到学习系统的端到端周期结构理论构建。
\end{remark}

\section*{本章七主题对照表}

\begin{center}
\begin{tabular}{>{\centering\arraybackslash}p{2.2cm}>{\centering\arraybackslash}p{10.5cm}}
\toprule
主题 & 本章对应 \\
\midrule
平移 & 梯度下降的参数平移与特征方向误差迭代 $e_t^{(i)} = (1-\eta\lambda_i)^t e_0^{(i)}$（\cref{ch32_gd_eigen}）；注意力窗口在序列上的滑动（\cref{def:self_attention_periodic}）；位置编码的平移不变性（\cref{rem:positional_encoding_translation_invariance}）；词嵌入的平移类比关系（\cref{prop:embedding_rotation_analogy}）；TransE 关系嵌入的平移结构（\cref{thm:relation_translation_rotation}） \\
缩放 & 学习率调度器的标度衰减（\cref{def:lr_scheduler}）；探索强度 $\varepsilon(t)$ 的指数衰减（\cref{def:exploration_exploitation}）；多头注意力的标度分解；词嵌入的模长向量（\cref{lem:embedding_phase_representation}） \\
旋转 & 动量法的阻尼旋转（\cref{prop:momentum_rotation}）；位置编码的复指数相位 $e^{i\omega\cdot pos}$（\cref{prop:positional_encoding_fourier}）；LSTM 门控的相位开闭（\cref{prop:lstm_long_range}）；词嵌入的旋转类比关系（\cref{prop:embedding_rotation_analogy}）；RotatE 关系嵌入的相位旋转（\cref{thm:relation_translation_rotation}）；推理链的相位累积（\cref{prop:reasoning_phase_accumulation}） \\
平衡 & 参数约束（权重衰减）、梯度传播约束（残差连接）、注意力归一化约束（softmax 归一化）（\cref{thm:agent_conservation}）；推理链的信息-深度权衡（\cref{lem:reasoning_depth_information}）；泛化-吸引子宽度权衡（\cref{prop:generalisation_attractor_width}） \\
对称 & 残差连接的恒等对称 $\mathbf{I} + \partial\mathcal{F}/\partial\mathbf{x}_l$；位置编码的平移不变相似度（\cref{rem:positional_encoding_translation_invariance}）；词嵌入分量的差角余弦对称（\cref{thm:embedding_cosine_structure}） \\
闭合 & 策略迭代向最优策略 $\pi^*$ 的有限步收敛闭合（\cref{prop:policy_iteration_cycle}）；多智能体参数平均的同步闭合（\cref{prop:parameter_averaging}）；推理链的迭代闭合（\cref{thm:reasoning_fixed_point}）；CoT 推理的相位锁定闭合（\cref{thm:cot_phase_locking}）；知识路径的相位闭合（\cref{prop:knowledge_path_period}） \\
周期 & 梯度下降的阻尼振荡（\cref{def:gradient_descent_oscillation}）；余弦退火周期 $2T$（\cref{prop:cosine_annealing_period_agent}）；多智能体同步周期 $T_{\text{sync}}$（\cref{prop:parameter_averaging}）；LSTM 门控的记忆周期（\cref{prop:lstm_long_range}）；Hopfield 极限环吸引子周期 $T_{\text{lc}}$（\cref{thm:hopfield_limit_cycle}、\cref{prop:limit_cycle_period}）；词嵌入的周期性语义生成（\textbf{猜想／经验假设}，注 \cref{prop:embedding_periodicity_generation}） \\
\bottomrule
\end{tabular}
\end{center}

\section*{本章思考题}

\begin{enumerate}[leftmargin=3em, label=\textbf{题 \arabic*}]
    \item 证明本章的\textbf{KDI 对周转圆参数的下降性}（\ref{prop:kdi_epicycle_descent_agent}）。提示：先写出本章目标的判定映射与演化群，再套用「周期系统」章通用判据的证明结构。

    \item 验证\textbf{智能体系统的三种约束关系}（\ref{thm:agent_conservation}）的前提在本章领域对象上成立，并给出一个不满足前提的反例。

    \item 根据本章七主题对照表，分别指出本章在\textbf{平移}、\textbf{缩放}、\textbf{旋转}三个主题上的具体载体，并各举一个领域数据实例。

    \item 综合题：用「周期系统」章的五步标准接口（表 \ref{tab:ps_five_steps}）重述本章的诊断流程，并逐一注明每一步所依赖的本章定理：\ref{prop:kdi_epicycle_descent_agent}、\ref{thm:agent_conservation}。
\end{enumerate}


